\chapter{Conclusioni}

\section{Sintesi del caso studio}

Il presente lavoro di tesi ha avuto come obiettivo la progettazione, l’implementazione e la valutazione di un sistema di \emph{Anomaly Detection} basato sull’algoritmo Isolation Forest, applicato all’analisi di log HTTP provenienti da un contesto reale. L’intero percorso di sviluppo è stato impostato con una finalità duplice: da un lato, verificare la fattibilità di un approccio automatico al rilevamento di comportamenti anomali nei log web; dall’altro, valutare in modo critico fino a che punto un modello non supervisionato possa essere utile in uno scenario applicativo concreto, caratterizzato da un forte squilibrio tra traffico lecito e richieste potenzialmente malevole.

Il punto di partenza dell’elaborato risiede nella constatazione che i sistemi di difesa tradizionali, fondati su regole statiche, firme note o blacklist manuali, mostrano limiti significativi nel fronteggiare la crescente dinamicità delle minacce informatiche e la variabilità del traffico web moderno. In un contesto di questo tipo, l’adozione di tecniche di Machine Learning rappresenta un’alternativa interessante, poiché consente di intercettare deviazioni dal comportamento atteso anche in assenza di una conoscenza preventiva di tutte le tipologie di attacco. Tale impostazione è particolarmente rilevante nei sistemi di analisi dei log, dove l’obiettivo non è soltanto individuare eventi noti, ma anche evidenziare pattern rari, sospetti o non immediatamente riconducibili a regole predefinite.

L’attività svolta si è concentrata sulla costruzione di un dataset derivato da log HTTP, sottoposto a un processo di pre-processing e di \emph{feature engineering} finalizzato a trasformare l’informazione testuale in variabili numeriche e booleane interpretabili dall’algoritmo. Le feature selezionate sono state progettate con un criterio orientato alla sicurezza applicativa: alcune descrivono caratteristiche dimensionali delle richieste, altre cercano pattern sintattici riconducibili a vettori d’attacco noti, altre ancora cercano di modellare aspetti comportamentali del traffico. Questo passaggio metodologico è stato centrale, perché nel caso di log web la qualità delle feature incide in modo determinante sulla capacità del modello di separare eventi normali, anomalie statistiche e vere anomalie di sicurezza.

L’algoritmo Isolation Forest è stato applicato dapprima a un sottoinsieme ridotto del dataset, così da consentire una prima fase esplorativa e una calibrazione ragionata dei parametri. In questo scenario limitato, il modello ha mostrato una discreta capacità di individuare gli attacchi inseriti manualmente, confermando che la combinazione tra feature engineering e anomaly scoring può produrre risultati utili quando il problema è ben circoscritto e il segnale anomalo è ancora distinguibile dal rumore di fondo. Questa fase ha quindi avuto un valore soprattutto dimostrativo e metodologico, oltre che sperimentale.

Tuttavia, l’estensione dell’analisi all’intero dataset ha mostrato con chiarezza i limiti dell’approccio adottato. Quando il modello viene applicato a un corpus di dimensioni molto elevate, con milioni di record e una quota di anomalie estremamente ridotta, la capacità discriminante dell’Isolation Forest si indebolisce sensibilmente. In particolare, i quattro attacchi inseriti manualmente non riescono più a emergere con affidabilità sufficiente rispetto alla massa dei log leciti, e il modello finisce per segnalare soprattutto osservazioni statisticamente atipiche ma non necessariamente malevole. Questo comportamento evidenzia un aspetto cruciale: in un contesto molto sbilanciato, il confine tra anomalia statistica e anomalia di sicurezza tende a sfumare, e il modello può privilegiare la forma del dato rispetto al suo significato operativo.

Un ulteriore elemento emerso dall’analisi riguarda il rapporto tra richiamo e precisione. L’Isolation Forest, soprattutto in configurazioni orientate alla massima sensibilità, tende a intercettare un numero elevato di outlier, ma questo si accompagna inevitabilmente a un aumento dei falsi positivi. Nel dominio della cybersecurity ciò non rappresenta soltanto un dettaglio tecnico, ma una questione sostanziale: un sistema che produce troppe segnalazioni non rilevanti rischia di ridurre la fiducia degli analisti e di abbassare l’efficacia operativa complessiva della soluzione. Di conseguenza, un buon modello di anomaly detection non deve limitarsi a “vedere tanto”, ma deve anche saper discriminare in modo sufficientemente fine ciò che merita realmente attenzione.

In questo senso, il lavoro conferma che l’Isolation Forest è uno strumento molto utile come base esplorativa, soprattutto quando si dispone di pochi esempi etichettati e si desidera ottenere una prima fotografia del comportamento del traffico. Il suo principale vantaggio è la scalabilità concettuale e computazionale; il suo principale limite, invece, emerge proprio quando il dataset cresce molto e il rumore supera il segnale. In tali condizioni, il modello tende a identificare soprattutto deviazioni statistiche generiche, mentre gli attacchi reali, se troppo pochi e troppo immersi nella distribuzione complessiva, possono perdere visibilità.

La lettura congiunta della matrice di confusione, del classification report e delle combinazioni di feature più frequenti tra le anomalie individuate porta a una conclusione chiara: il sistema non va interpretato come un sostituto dei meccanismi di difesa tradizionali, ma come uno strumento complementare di supporto all’analisi. In un contesto operativo reale, un approccio di questo tipo dovrebbe essere integrato con ulteriori filtri, controlli a valle, regole di priorizzazione o componenti supervisionate, così da ridurre l’impatto dei falsi positivi e aumentare la probabilità di intercettare gli eventi davvero critici.

\section{Considerazioni finali}

Dal punto di vista metodologico, l’esperienza svolta mette in evidenza alcuni elementi di valore generale. In primo luogo, mostra che l’analisi dei log HTTP può fornire informazioni estremamente utili non solo per la diagnostica e il monitoraggio, ma anche per attività di sicurezza orientate al rilevamento automatico di pattern anomali. In secondo luogo, evidenzia che il successo di un progetto di anomaly detection dipende in larga misura dalla qualità della rappresentazione dei dati: scegliere le feature giuste, definire correttamente le etichette e interpretare in modo critico gli score prodotti dal modello sono passaggi decisivi quanto la scelta dell’algoritmo stesso.

Il caso studio ha inoltre mostrato che, nei problemi di sicurezza applicata, le metriche vanno lette con attenzione e mai isolate dal contesto. Un risultato apparentemente alto in termini di accuracy può risultare fuorviante se il dataset è fortemente sbilanciato; allo stesso modo, una buona recall sugli attacchi non garantisce automaticamente l’utilizzabilità del sistema, se la precisione resta troppo bassa e il numero di falsi positivi diventa eccessivo. In altre parole, non basta che il modello sia “tecnicamente corretto”: deve essere anche coerente con l’obiettivo operativo per cui viene impiegato.

Un altro aspetto importante riguarda il ruolo della conoscenza di dominio. Nel caso dei log HTTP, l’interpretazione dei risultati beneficia fortemente della comprensione del contesto applicativo: richieste lunghe, crawler, strumenti automatici di scansione e pattern sintattici sospetti non hanno tutti lo stesso significato. Alcuni rappresentano segnali potenzialmente malevoli, altri sono semplicemente caratteristiche atipiche di traffico legittimo. La capacità di distinguere questi casi richiede non solo un modello, ma anche un’analisi critica umana e, idealmente, un ciclo di feedback continuo tra sistema automatico e validazione esperta.

Sotto il profilo dei risultati, il lavoro conferma quindi un punto fondamentale: l’Isolation Forest è efficace come **strumento esplorativo**, ma non è sufficiente, da solo, a risolvere il problema del rilevamento automatico di anomalie in contesti molto grandi e rumorosi. Nel sottoinsieme ridotto del dataset il modello riesce a evidenziare gli attacchi inseriti manualmente con una certa chiarezza; sull’intero corpus, invece, la sua capacità di generalizzazione si riduce sensibilmente, e gli attacchi finiscono per essere assorbiti dalla massa dei log normali e dalle anomalie statistiche non rilevanti dal punto di vista della sicurezza. Questo non invalida il metodo, ma ne delimita con precisione il campo di applicazione.

In conclusione, il lavoro svolto rappresenta un passo concreto verso sistemi di analisi dei log più intelligenti e adattivi, ma evidenzia anche che l’automazione del rilevamento non può prescindere da un’integrazione con altre forme di controllo. Le possibili evoluzioni future includono l’uso di modelli ibridi, l’introduzione di meccanismi supervisionati o semi-supervisionati, l’affinamento delle feature, l’ottimizzazione della soglia decisionale e l’adozione di strategie di post-processing basate sul feedback dell’analista. L’obiettivo non dovrebbe essere soltanto ridurre il numero di errori, ma costruire un sistema capace di trasformare i segnali grezzi del traffico web in informazioni davvero utili, affidabili e azionabili.